\section{引言}
\label{sec:intro}

\subsection{基因编辑工具设计中的预测与理解}

CRISPR-Cas9 通过可编程的 RNA--DNA 配对实现定点切割，sgRNA 的活性差异使"实验前筛选高效位点"成为核心计算问题\citep{Jinek2012Science,Cong2013Science}。早期方法依赖序列组成与位置偏好的人工规则\citep{Doench2014NatBiotechnol,MorenoMateos2015NatMethods}，随后被数据驱动的模型取代：DeepBind 与 Basset 证明卷积网络可从序列实验数据中学习局部序列特异性\citep{Alipanahi2015NatBiotechnol,Kelley2016GenomeRes}；DeepCRISPR 进一步把序列与细胞特异性表观遗传信息（CTCF、DNase、H3K4me3、RRBS）纳入统一编码，用于 sgRNA 活性与脱靶预测\citep{Chuai2018GenomeBiol}；Rule Set 2 等方法也显著提升了预测能力\citep{Doench2016NatBiotechnol}。

\subsection{仅靠预测无法回答影响因素问题}

预测能力的提升并不意味着研究者能够理解模型依据哪些因素作出判断。不同模型类别的可解释性各有边界：线性模型给出方向与统计量但表达受限；树模型可建模非线性与交互\citep{Chen2016KDD}，但 Gain/Weight/Cover 没有方向；深度模型的内部表示难以直接解释。SHAP 与 Integrated Gradients 分别从 Shapley 值与公理化归因出发提供了通用解释工具\citep{Lundberg2017NeurIPS,Sundararajan2017ICML}，Transformer 则以自注意力建模跨位置关系\citep{Vaswani2017NeurIPS}。然而，即使归因可用，仍需回答：哪些序列位置与模式与效率相关？表观遗传环境是否在序列之外提供**额外预测信息**？同一因素在不同细胞背景下是否稳定？不同模型是否**独立**支持同一候选因素？

\subsection{从"模型解释"到"科学发现流程"}

现有研究中，预测、归因、统计与跨条件比较往往分散在不同脚本与不同报告里，导致结论难以逐项追溯，也容易把模型依赖强度误读为生物学效应。本研究的目标是构建一个**科学发现平台**：把上述分析组织为统一、可复现、可追溯的流程，使每一步结论都能回溯到具体结果文件，并明确区分：预测效应（$\Delta R^{2}$）、模型归因（attribution）、稳健性（如信噪比）与统计证据（$p$/FDR）。

平台围绕三个研究问题组织分析：（1）**序列**：哪些位置与候选模式与编辑效率相关？（2）**环境**：细胞表观遗传信息是否提供额外的\emph{预测}信息？（3）**背景**：同一因素在不同 cell line 中是否保持一致？在此框架中，预测模型是发现影响因素的工具，而不是最终结论。

\subsection{案例研究与贡献定位}

本文以 DeepCRISPR 公开数据作为**案例研究（case study）**，检验平台能否从多模型、多条件的受控实验中发现可解释、可追踪、可比较的候选因素。案例研究**不**用于主张平台在预测精度上优于已有方法，也**不**用于证明表观遗传环境必然提升预测——后者的检验结果在本案例中为弱证据（见 \S\ref{sec:res-environment}）。平台的贡献在于：（i）统一的、以结果文件为接口的多模型发现流程；（ii）在该流程中同时保留稳定证据与模型特异性/上下文依赖证据；（iii）从计算证据出发给出可实验检验的候选假设与最小扰动方案。
